% GENERATED from QPw00-page-loop.md by md2tex.py. Do not hand-edit: sync is
% one-way and the next run overwrites this file.

\section{DRAFT gives the Page a promise}
\textbf{The DRAFT contract}: what enters, what DRAFT may decide, and what it hands forward.

\begin{verbatim}
 DRAFT
+--  enters    page need . constraints . prior round
+--  owns      purpose . Aims . promised shape
+--  writes    any Page part needed to expose the design
+--  names     unknowns . dependencies . assumptions
+--  exits     one stable round contract
\end{verbatim}
📌 DRAFT decides what this round of the Page is trying to become.

DRAFT's full contract, its worked example turning one approved Point into several sentence scaffolds, and its rule that the Aim behind a hole is the stake a probe card will later carry all live at QPw2-draft. OUTLINE, at QPw1-outline, is the phase now immediately upstream of DRAFT: it agrees the plan's shape before DRAFT turns an approved bullet into prose, a split this page's own §1 predates and never had to name.


\section{PROBE resolves what the Page cannot know}
\textbf{The PROBE contract}: an approved mark leaves the plan, becomes a card, and is asked.

\begin{verbatim}
 PROBE
+--  enters    the approved mark, still bare, from the plan
+--  owns      the MATCH lookup, the card, and the dispatch
+--  writes    probe/PP<NN>-<slug>/ . consumer/ . executor/
+--  exits     the moment the stripped question leaves
+--  avoids    target Page prose . landing the answer
\end{verbatim}
📌 PROBE turns an approved mark into a card and asks; landing what comes back is EVIDENCE's job.

PROBE's trigger, its four-step MATCH lookup that runs before any dispatch, and its optional-when-no-unknown-is-hidden rule now live at QPw3-probe. The write surface it hands off to, binding the returned answer and landing it in a card, now lives at QPw4-evidence. PROBE's meaning changed on 260817, exactly as this page's own Opening states it: it now only raises the card and asks.

The target Page is the Page whose lifecycle raised the stake-bearing Q-consumer. A family may file the QA-probe under an evidence route such as Literature or Value without transferring the Q-consumer, A-consumer, or State to that route's topic Page. The Paper family's S03 and S04 layout exposes the distinction: a Results Page can raise the question, a QA-probe can live beneath a Value topic, and a QA-bank file can answer it. Treating physical placement as an ownership transfer gives one exchange two Page-facing consumer surfaces before the answer reaches prose.

One PROBE run may write its declared Probe surface and the active target Page's Probe reference, A-consumer, and State. When the same Q-executor serves other Q-consumers, the Probe surface records their references and makes the answer available without authoring those sibling Pages in the current run. Each sibling interprets the answer and changes its own Content through its own PROBE or REVISE route, with its own version and CHECK. This preserves one Q-executor for reuse without turning one answer into an unbounded cross-Page edit.

The active Page may need to show the Q-executor, bank target, state, returned answer, and limits while keeping its source concise. A read-only projection can render that chain from the QA-probe and bank answer, while the Page source keeps the Probe reference and its own A-consumer interpretation. The same reference should drive phase-scoped context loading and CHECK, so an author does not maintain a Probe pointer, an embed, and a Related Board Pages row for one relationship. Missing targets, superseded answers, or changed source hashes must fail visibly rather than leave an old projection looking current.


\section{REVISE makes the current promise work}
\textbf{The REVISE contract}: the Page changes while its purpose and Aims remain fixed.

\begin{verbatim}
 REVISE
+--  enters    current Page . evidence . feedback
+--  holds     purpose . Aims
+--  owns      add . delete . move . rewrite
+--  updates   Content . Opening . States . records
+--  exits     stronger version . visible remaining gaps
\end{verbatim}
📌 REVISE improves the realization of the current round rather than redefining its promise.

REVISE's fixed-promise test, its five-step LAND, ARGUE, SOUND, RENDER, BUILD order, and the three display-walk steps it owns now live at QPw5-revise, which also carries COMPILE. This page's own §6, Operations route by reason, still carries the shared DRAFT-versus-REVISE test table that separates the two phases on the same edit, so a reader is pointed sideways as well as forward.


\section{CHECK decides where the current version goes next}
\textbf{The CHECK contract}: one concrete version is judged against its promise and routed.

\begin{verbatim}
 CHECK
+--  enters    rendered version . Aims . evidence . constraints
+--  owns      judgment . findings . closing decision
+--  writes    comments . findings . gate record
+--  routes    close . REVISE . EVIDENCE . new DRAFT . HOLD
+--  avoids    curing its own substantive finding
\end{verbatim}
📌 CHECK observes and decides; another phase performs the content change it requests.

CHECK's judge-may-not-repair separation, the three independent built-artifact counts it reads (declared, rendered, accepted), and its accept-biased human gate now live at QPw6-check. This page's own §10 still carries the receipt fields and invariants CHECK's routed findings are measured against, so nothing here is left an orphaned reference.


\section{Transitions form rounds, not a rigid conveyor belt}
\textbf{The transition grammar}: repetition is legal, and returning to DRAFT changes the round.

\begin{verbatim}
 ROUND n . purpose + Aims fixed

 DRAFT  --+-->  CHECK --> [x] close
              +-->  PROBE  -->  REVISE  -->  CHECK
                                       |                 |
                        +---------------+                 |

 ROUND n+1 <--------------  DRAFT <---- purpose or Aims reopened
\end{verbatim}
📌 The arrows express dependencies and routing choices, not a rule that every phase runs once.

\texttt{DRAFT, DRAFT, PROBE, REVISE, DRAFT} is a legal history. The first two DRAFT entries can be repeated design work in one round. The final DRAFT means that revision reopened the promise and therefore began a new round. A complete draft with no unknowns and no revision need may go directly to CHECK.

REVISE may reveal that the current purpose or Aims are wrong. The Page then returns to DRAFT and receives a new round, while the Page itself persists. Earlier evidence remains available, but its relevance must be checked against the new promise. Any earlier closing decision no longer applies to the new round.

The history should say which authority was used: promise reopened, unknown resolved, current promise improved, or version judged. Without that reason, two identical diffs can be mislabeled as different phases and two different intentions can be mislabeled as the same one.


\section{Operations route by reason, not by edit shape}
\textbf{The same operation under two authorities}: a compact test for ordinary Page changes.

\begin{verbatim}
 operation    DRAFT when                  REVISE when
--------------------------------------------------------------------
 add         new promise or Aim             serves a fixed Aim
 delete      removes a promise or Aim        removes weak or extra prose
 move        changes the promised argument  improves flow under the promise
 rewrite     reframes purpose or Aim         improves supported expression
\end{verbatim}
📌 The operation says what changed in the file; the authority says which phase performed it.

Adding a paragraph to explain evidence for an existing Aim is REVISE. Adding a paragraph because the Page now answers a new question is DRAFT. Deleting unsupported or duplicate prose while keeping the Aim is REVISE. Deleting the promised result itself, or deleting its Aim, is DRAFT.

Moving paragraphs to improve flow under the same argument is REVISE. Changing the argument the Page promises to make is DRAFT even if the diff is only one moved heading. Rewriting the Opening for clarity is REVISE when the purpose stays fixed. Rewriting it to give the Page a different purpose is DRAFT.

PROBE writes questions, sources, evidence, and answer records. CHECK writes findings, comments, and gate records. When either phase causes target prose to change, the actual content edit is performed under REVISE or a restarted DRAFT.


\section{One lifecycle Page holds the phases together}
\textbf{The Page split rule}: each phase begins as a Content division and earns a Page only by becoming an independent question.

\begin{verbatim}
 QB5 . ONE LIFECYCLE QUESTION
+--  1 .  DRAFT
+--  2 .  PROBE
+--  3 .  REVISE
+--  4 .  CHECK
+--  5 .  transitions
+--  6 .  operations

 split test . ALL required
+--  independent question
+--  own Aims +  States
+--  independent closing gate
+--  own continuation files

 phase skill   executable contract
 design Page   independently closable question
 mapping       no automatic 1:1 mirror
\end{verbatim}
📌 Shared boundaries stay together until one phase can carry and close a question of its own.

DRAFT, PROBE, REVISE, and CHECK remain divisions of QB5 because their meanings depend on comparison and transition. The DRAFT and REVISE boundary is easier to judge when both definitions and the operation examples stay on one Page. The transition rules also need one home that no phase-specific Page can own alone.

The phase must have an unresolved question that is not merely asking for its definition. It must need its own Aims and States rather than borrowing QB5's records. It must have an independent closing gate and its own small continuation map in Files. If any test fails, the material remains a division or paragraph on QB5.

A phase skill may remain separate because a worker needs to load one execution contract at a time. That file boundary does not force the Board to create a matching design Page. If a phase later passes the split test, QB5 keeps the shared boundaries and transitions while pointing to the new Page for that phase's independent question.


\section{Page Type and Page Phase are separate skill axes}
\textbf{The skill composition}: the base resolves what the Page is and how the current work is acting on it.

\begin{verbatim}
 haipipe-page                   the shared Page contract and router
+--  page-types/                      what kind of Page persists . ten types . QB6 owns the roster
|   +-- for-stage
|   +-- for-skill
|   +-- for-venue
|   +-- ... seven more . for-design . for-display . for-literature . for-meeting . for-section . for-slide . for-value
+--  page-phases/                     what authority acts now
    +-- draft
    +-- probe
    +-- revise
    +-- check

one invocation = base + matching Page Type + current Page Phase + family worker
\end{verbatim}
📌 Page Type and Page Phase are orthogonal, so their folder names and skill names must not collapse them into one label.

\texttt{haipipe-page-for-stage}, \texttt{-for-skill}, and \texttt{-for-venue} keep their names and move under \texttt{page-types/}. The roster has since grown to ten Page Types, and \texttt{QB6} owns the admission test and the list. The grouping folder is organizational and carries no \texttt{SKILL.md} of its own. A new Page Type is added only when a persistent Page needs a structural contract that the base does not provide.

The phase contracts are \texttt{haipipe-page-draft}, \texttt{-probe}, \texttt{-revise}, and \texttt{-check} under \texttt{page-phases/}. They apply across Page Types and therefore do not use \texttt{for-stage} in their names. The base first adopted the phase vocabulary without adding \texttt{ADVANCE}. The automatic router now earns a verb named \texttt{RUN}, because it may repeat, branch, HOLD, or return to DRAFT rather than advance in one direction.

The target Page owns the stake-bearing Q-consumer. The PROBE phase strips the stake into a neutral Q-executor, binds the returned A-executor, and writes an A-consumer interpretation for each consumer. One Q-executor may serve several Q-consumers. \texttt{haipipe-probe} remains the shared crossing protocol, while \texttt{haipipe-page-probe} applies that protocol to a Board Page.


\section{RUN turns the phase grammar into a bounded loop}
\textbf{The executable flow}: one controller composes separate producer, builder, and judge roles without prescribing one phase sequence.

\begin{verbatim}
 raw-material packet
|  Page . Type . start Phase . intent . sources . constraints . gate . limits

 controller
+--  DRAFT /  PROBE /  REVISE  -> producer
+--  build + version snapshot       -> mechanical builder
+--  CHECK exact version            -> fresh read-only judge
                  |
                  +-- [x] CLOSE
                  +--  REVISE ------+
                  +--  PROBE -------+
                  +--  DRAFT round+1|
                  +--  HOLD          |
                                      +--> controller
\end{verbatim}
📌 RUN follows returned authority routes and stops at explicit terminals or limits. Figure~\ref{fig:qpw1-run-loop} states the same controller as an algorithm, so a reader can see the three actors alternate and read the stop conditions in the order the loop tests them.

\begin{figure}[H]
\centering
\includegraphics[width=.85\linewidth]{../display/QPw00-Display1-run-loop/assets/figure.pdf}
\caption{RUN, as the router it is rather than the conveyor belt its name could
  suggest. The producer, the builder, and the judge are three separate actors;
  the loop follows the route each returns, refuses a route outside $R(\varphi)$,
  and stops at CLOSE, HOLD, or a limit.}
\label{fig:qpw1-run-loop}
\end{figure}

ADVANCE suggests that one phase has one next phase and that progress always moves forward. RUN means execute the current Page lifecycle from a named starting authority until CLOSE or HOLD. The legal route table is dynamic: DRAFT, PROBE, and REVISE may repeat or hand off, while only CHECK may CLOSE. Figure~\ref{fig:qpw1-route-relation} writes that table as two rules and derives both laws from them, including one this prose never states: CHECK cannot route to CHECK, so a judged version reaches a second judgment only through a producer. Returning to DRAFT from another phase increments the round only when purpose or an Aim reopened.

\begin{figure}[H]
\centering
\includegraphics[width=.85\linewidth]{../display/QPw00-Display2-route-relation/assets/figure.pdf}
\caption{The legal-route relation $R$, seven rows, and the three laws the
  table makes checkable. The rows are deliberately NOT uniform: PROBE may not
  hand off to DRAFT or CHECK, and REVISE may not hand off to PROBE. The judge
  alone reaches CLOSE and alone cannot reach itself, so every judged version
  passes through a producer before it is judged again; and every phase without
  exception may stop at HOLD.}
\label{fig:qpw1-route-relation}
\end{figure}

The packet names the persistent Page, stable Page Type, starting Phase, run intent, source paths, settled constraints, closing gate, and step and round limits. A new Page is first created and registered, then RUN starts it at DRAFT. An existing Page with no known next need starts at CHECK, allowing a cold judge to route the visible version. A missing source, unknown gate, or ambiguous authority becomes HOLD rather than an invented input.

The producer performs exactly one DRAFT, PROBE, or REVISE phase and suggests a legal route. The mechanical builder rebuilds, runs deterministic checks, and identifies the source plus render version. The fresh reviewer performs CHECK against that exact version and returns CLOSE, REVISE, PROBE, DRAFT, or HOLD. The controller validates and follows the route, but it cannot alter Page prose or convert a pending human gate into CLOSE.

The run stops on CLOSE, explicit HOLD, missing input, failed build, version mismatch, required human ruling, maximum steps, or maximum rounds. Reaching a limit says the process did not converge within its budget. It never says the Page passed. The run's final state and every attempted phase remain inspectable instead of disappearing into an agent transcript.


\section{Audit proves process claims with receipts and fault tests}
\textbf{The assurance model}: deterministic invariants, fresh semantic judgment, and direct or human evidence support different quality claims.

\begin{verbatim}
 phase receipts             what happened, by whom, to which version
 deterministic audit       legal routes . rounds . roles . versions . bounds
 fresh CHECK               function . evidence . readability . local gate
 human evidence            approval only where the Page Type requires it
 branch + fault tests      expected success and expected refusal
\end{verbatim}
📌 A passing process audit proves that the declared process ran correctly, not that every possible substantive claim is true.

Each receipt records step, round, Phase, producer or judge actor, builder actor, status, source and render SHA-256 values, route, reason, artifacts, evidence, findings, and human-gate state. CHECK additionally binds \texttt{checked\_version} and a verdict. The receipts are ordered and stored outside Page discovery under \texttt{\_runs/page/<page-id>/<run-id>.json}. The deterministic auditor independently rehashes the source and rendered files on disk; matching claims repeated inside a receipt are not accepted as artifact evidence by themselves. The terminal CHECK record is not appended to the Page afterward, because that append would change the version it claims to approve.

The preserved raw-material packet must match the run identity, first Phase, declared gate, and limits. Only legal phase routes are accepted, and only CHECK may CLOSE. A non-DRAFT route to DRAFT must name the reopened purpose or Aim and increment the next round exactly once. The producer, mechanical builder, and judge of one version must have different actor identities. Every version id must be the two declared lowercase SHA-256 digests joined by \texttt{:}, and every receipt must begin from the preceding receipt's ending version. CHECK must observe identical before, after, and checked ids; the auditor rehashes the current artifacts, and any later change requires a new CHECK. A required human gate closes only with durable evidence that the person ruled. Maximum steps and rounds terminate as non-convergence, never as a pass.

Happy paths include DRAFT directly to CHECK and DRAFT through optional PROBE and REVISE. Branch tests include CHECK to REVISE and back, CHECK to PROBE, and CHECK to a new DRAFT round. Fault injection includes self-approval, mutation after CHECK, illegal route, packet mismatch, broken version continuity, symbolic hashes, missing human evidence, failed worker, non-terminal trace, and exhausted limits. The deterministic harness must reject each injected fault for the specific invariant it violates.

The mechanical checker can prove structural facts, and the lifecycle auditor can prove process facts. A fresh reviewer supplies semantic evidence that the Page performs its declared function and is readable without the drafting conversation. Direct sources or a human gate supply claims those instruments cannot settle. The final report therefore names the checked version, traversed branches, evidence inspected, remaining findings, gate state, and residual risk instead of saying only that quality is guaranteed.


\section{Workflow, phase, step and round are four different words}
\textbf{The vocabulary rule}: the loop needs all four and none of them substitutes for another, which is why "phase" was kept when JL weighed replacing it with "step" on 260818.

\begin{verbatim}
word         answers                      count in the ONE live run    repeats?
----------------------------------------------------------------------------------
 WORKFLOW  which LOOP is this?          1  . the run itself          no
 PHASE     which AUTHORITY acts?        7 defined . 2 used           YES
 STEP      WHERE in this run?           5 . monotonic 1..5           never
 ROUND     which PROMISE era?           1 . bumps only on a reopen   on reopen
\end{verbatim}
📌 A phase is a TYPE and a step is an INSTANCE of one: in \texttt{260805-0216-QB8e} the single CHECK phase occupies steps 1, 3 and 5.

Renaming phase to step would put two different meanings on one key inside a single receipt, and the auditor reads both. \texttt{step} is the monotonic position that never recurs; \texttt{phase} is the kind of work that must be free to recur, which is the property the whole loop is built on.

The word must permit repetition, because CHECK may route back to any earlier phase, PROBE may be skipped entirely, and a page may re-enter DRAFT in a new round. The one live run proves it rather than asserting it: CHECK ran three times and REVISE twice inside a single round, so a vocabulary that numbered them 1..7 could not describe what happened. JL raised the replacement and answered it himself on 260818: "for phase, we can do it again and again, right?"

The directory says which of the four words scopes the family, and the contracts inside it say which authority each phase holds. So the full term is "workflow phase" where disambiguation is needed and "phase" everywhere else, and the 927 occurrences of the word across this board and the skill tree stay as they are.


\section{The same loop, priced: what a person is actually asked to do}
\textbf{The cost rule}: a reader deciding whether to run this loop needs the WORK, not the authority, so the loop is stated once more as one row per step, with a person's job and its price in it.

\begin{verbatim}
#    PHASE      WHAT HAPPENS                    YOUR JOB              TIME    HOW OFTEN
-------------------------------------------------------------- once per page -----------
0    OPEN       you name the page and the        say what it        15 min  once
                one thing it must decide           must decide
------------------- the loop . CHECK may send you back to any row in this band ---------
1    OUTLINE    a machine writes the SHAPE       approve it, or     10 min   each round
                only: sections, bullets, and       send it back
                what each bullet still owes      reject costs 10 sec
2    DRAFT      a machine turns each approved   nothing.              0        each round
                bullet into sentences with      do NOT fill a hole
                visible <HOLE>s
3    PROBE      a machine looks for an answer   nothing               0        each round
                that already exists, then opens
                a card and asks the bank
4c   CITATION   a PERSON finds the published     paste the entry,   20 min   per entry
                work; a machine may copy           tick `verified`
                bibtex and never compose it
4v   VALUE      the bank answers, bound BY       read the proof,    30 min   per card
                PATH to a real file                tick `read:`
                                                 the cost, and the tick that permits quoting
4d   DISPLAY    a machine freezes the exact     nothing               0        per figure
                bytes a figure is drawn from
                 4c . 4v . 4d run at the SAME TIME. None waits for another to finish.
5    REVISE     a machine replaces every landed nothing               0        each round
                hole, draws the figures, and
                rebuilds the pdf
6    CHECK      a machine and a DIFFERENT        accept, or say     15 min   each round
                reader judge the BUILT pdf         where to go back
                 "go back"  row 1, 2, 3, 4 or 5
------------------------------ reached only when CHECK says CLOSE ----------------------
7    ACCEPT     each figure is bound to the      tick               10 min  once per figure
                inputs it was accepted with        `accepted: [x]`
                                                 reverts if the inputs change
8    CLOSE      the page's own ruling, named     sign it            5 min   once per round
                by its Page Type                 this ends the round
\end{verbatim}
📌 SEVEN of the eleven rows ask something of a person and four read \texttt{nothing}; of those seven, five carry a tick a machine may never write, and two (row 0 OPEN and row 6 CHECK) ask for a judgment that leaves no tick behind. That is the same fact \texttt{§9} states as authority, priced instead. Figure~\ref{fig:qpw00-who-does-what} renders this table for print, with the eight files it was transcribed from frozen beside it.

\begin{figure}[H]
\centering
\includegraphics[width=.85\linewidth]{../display/QPw00-Display3-who-does-what/assets/figure.pdf}
\caption{The page workflow as a person experiences it: one row per step, what
  the machine does, what a person must do, and how often. SEVEN of the eleven
  rows ask something of a person and four read \emph{nothing}; of those seven,
  five carry a tick a machine may never write and two (OPEN and CHECK) ask for
  a judgment that leaves no tick behind. TIME is the only estimated column,
  because no OUTLINE, DRAFT, PROBE or EVIDENCE phase has ever been timed on
  this board.}
\label{fig:qpw00-who-does-what}
\end{figure}

No OUTLINE, DRAFT, PROBE or EVIDENCE receipt exists anywhere on this board, so four of the table's phases have never run under the contract at all. Every other cell is transcribed from a file frozen in \texttt{display/QPw00-Display3-who-does-what/intake/inputs/}, and the manifest names all eight with their sha256. A number nobody has measured is marked as an estimate on the table's own face rather than left to look like the others.

A reader who takes only the three bands away already has the loop: one thing happens once, six things repeat until CHECK stops sending them back, and two things happen when it does. The numbers are not a sequence, because rows 4c, 4v and 4d run at the same time and row 6 may return to any of rows 1 through 5. That is the non-linearity \texttt{§5} argues from transitions, put in the one place a person looks to find out what today costs them.

